% 在本目录运行 latexmk -xelatex example.tex
\documentclass[a4paper,12pt]{article}
\usepackage{buaa-proposal}
\usepackage{amsmath,booktabs}
\newcommand{\ReportTitle}{面向复杂应用场景的信息管理系统设计与实现及性能优化研究}
\newcommand{\StudentName}{示例姓名}
\newcommand{\StudentID}{00000000}
\newcommand{\SchoolSupervisor}{指导教师}
\newcommand{\CompanySupervisor}{}
\newcommand{\SupervisorCompany}{}
\newcommand{\ReportYear}{2026}
\newcommand{\ReportMonth}{11}
\newcommand{\ReportDay}{18}
\begin{document}
\hypersetup{pageanchor=false}
\MakeCover
\ReportDecorationstrue
\pagenumbering{Roman}
\hypersetup{pageanchor=true}
\tableofcontents
\pagenumbering{arabic}
\section{研究背景与目标}
本示例用于检查长题目封面，以及正文、公式、图表和列表之间的排版。图\ref{fig:logo}和表\ref{tab:plan}演示交叉引用。相关方法见文献\cite{sample-a, sample-b}，重复引用见\cite{sample-a}。

这是直接输入的普通段落。正文采用十二磅字体和固定的基本行距。较高的公式会撑开所在行，例如 $\displaystyle\frac{\sum_{i=1}^{n}x_i^2}{1+\sqrt{x}}$。

这是另一个普通段落，用空行分段。
\subsection{公式与列表}
独立公式放在段落命令外部，可以使用标准公式环境。

\begin{equation}
\label{eq:cost}
C=\sum_{i=1}^{n}\frac{w_i}{1+x_i}.
\end{equation}
式\eqref{eq:cost}给出示意目标函数。以下列表列出验证步骤。

\begin{enumerate}
\item 收集实验数据。
\item 比较不同参数配置。
  \begin{itemize}
  \item 记录内存占用。
  \item 统计处理时间。
  \end{itemize}
\end{enumerate}
列表结束后继续正文。

\begin{figure}[htbp]
\centering
\includegraphics[width=.42\linewidth]{assets/school-logo.jpeg}
\caption{示例图片}\label{fig:logo}
\end{figure}
\begin{table}[htbp]
\centering
\caption{工作计划}\label{tab:plan}
\begin{tabular}{ll}
\toprule
阶段 & 工作内容 \\
\midrule
第一阶段 & 文献阅读与需求分析 \\
第二阶段 & 实现与实验验证 \\
\bottomrule
\end{tabular}
\end{table}
\section{长段落与分页}
\subsection{面向不同任务规模与资源约束的系统设计、性能评估与实验结果分析方法}
\subsubsection{连续正文}
\newcount\examplecount
\examplecount=0
\loop
这一段用于检验连续正文的自然分页。研究将结合任务规模、资源约束和实际运行情况开展实验，记录系统吞吐量、响应时间和内存占用，并讨论不同配置对实验结果的影响。

\advance\examplecount by 1
\ifnum\examplecount<18
\repeat
\begin{thebibliography}{99}
\bibitem{sample-a} 示例作者. 示例文献题名. 2026.
\bibitem{sample-b} 另一位示例作者. 实验设计方法. 2026.
\end{thebibliography}
\end{document}
